% =============================================================================
%                    WEEK 2: TRANSFORMATIONS & KINEMATICS
% =============================================================================
\section{Week 2: Representations, Transformations \& Kinematics}

\subsection{Topic Summary}

\begin{keybox}[Learning Objectives]
By the end of this week, you should be able to:
\begin{itemize}
  \item Apply 2D and 3D rigid body transformations (translation, rotation).
  \item Construct and multiply homogeneous transformation matrices.
  \item Explain Euler angles, gimbal lock, and quaternions for 3D rotations.
  \item Compute forward kinematics for kinematic chains.
  \item Set up and solve inverse kinematics problems analytically.
  \item Explain the role of the Jacobian in numerical inverse kinematics.
\end{itemize}
\end{keybox}

\textbf{Big Picture.} To plan robot motion, we need to mathematically describe \emph{where} things are. This week covers rigid body transformations---the mathematical tools for specifying position and orientation in 2D and 3D. We then extend to kinematic chains (robot arms) and the fundamental problems of forward and inverse kinematics.

\textbf{What Examiners Like to Ask:}
\begin{itemize}
  \item ``Compute the homogeneous transformation matrix for rotation $\theta$ then translation $t$.''
  \item ``Apply a sequence of transformations to find the final pose.''
  \item ``Explain gimbal lock and how quaternions avoid it.''
  \item ``Derive forward kinematics for a 2-link planar arm.''
  \item ``Solve inverse kinematics for a 2-link arm given end-effector position.''
  \item ``What is the Jacobian and how is it used in IK?''
\end{itemize}

\bigseparator

\subsection{Core Concepts}

\textbf{Roadmap:}
\begin{enumerate}
  \item Rigid body transformations in 2D
  \item Homogeneous coordinates and transformation matrices
  \item 3D rotations: Euler angles, axis-angle, quaternions
  \item Kinematic chains and forward kinematics
  \item Inverse kinematics
\end{enumerate}

% =============================================================================
%                         2D TRANSFORMATIONS
% =============================================================================
\subsubsection{Rigid Body Transformations}

\begin{defbox}[Rigid Body Transformation]
A \term{rigid body transformation} $h: A \to W$ maps every point of body $A$ into workspace $W$ while:
\begin{enumerate}
  \item Preserving distances between all pairs of points.
  \item Preserving orientation (no mirror images).
\end{enumerate}
\end{defbox}

A rigid body in 2D is fully described by:
\begin{itemize}
  \item \textbf{Position:} $(x, y) \in \mathbb{R}^2$
  \item \textbf{Orientation:} angle $\theta \in [0, 2\pi)$
\end{itemize}

\separator

\subsubsection{2D Translation}

\begin{defbox}[Translation in 2D]
Translation by $(x_t, y_t)$:
\[
h(x, y) = (x + x_t, y + y_t)
\]
For primitives $H_i = \{(x,y) : f(x,y) \leq 0\}$:
\[
h(H_i) = \{(x,y) \in W : f(x - x_t, y - y_t) \leq 0\}
\]
\end{defbox}

\separator

\subsubsection{2D Rotation}

\begin{defbox}[Rotation Matrix in 2D]
Rotation about the origin by angle $\theta$ (counter-clockwise):
\[
R(\theta) = \begin{pmatrix} \cos\theta & -\sin\theta \\ \sin\theta & \cos\theta \end{pmatrix}
\]
A point $v = (x, y)^T$ transforms to:
\[
v' = R(\theta) \cdot v
\]
\end{defbox}

\begin{thmbox}[Properties of 2D Rotation Matrices]
\begin{itemize}
  \item \textbf{Orthogonal:} $R^T R = R R^T = I$
  \item \textbf{Determinant:} $\det(R) = 1$
  \item \textbf{Inverse:} $R(\theta)^{-1} = R(-\theta) = R(\theta)^T$
  \item \textbf{Composition:} $R(\theta_2) R(\theta_1) = R(\theta_1 + \theta_2)$
  \item \textbf{2D rotations commute:} $R(\theta_1) R(\theta_2) = R(\theta_2) R(\theta_1)$
\end{itemize}
\end{thmbox}

\begin{tipbox}[Right-Hand Rule]
Positive rotation is \textbf{counter-clockwise} (direction fingers curl when thumb points along rotation axis). In 2D, rotations are about the $z$-axis extending ``out of the page.''
\end{tipbox}

\separator

\subsubsection{Complex Numbers for 2D Rotations}

\begin{defbox}[Rotation via Complex Multiplication]
Complex numbers with unit norm can represent 2D rotations:
\[
z = \cos\theta + i\sin\theta = e^{i\theta}
\]
Rotating point $(x, y)$ by $\theta$: multiply complex number $x + iy$ by $e^{i\theta}$.
\end{defbox}

\separator

\subsubsection{Homogeneous Coordinates}

\begin{defbox}[Homogeneous Coordinates]
A 2D point $(x, y)$ is represented in \term{homogeneous coordinates} as $(x, y, 1)^T$.

More generally, $(x, y, w)$ represents the Cartesian point $(x/w, y/w)$ for $w \neq 0$.

Points at infinity: $(x, y, 0)$ represents infinity in direction $(x, y)$.
\end{defbox}

\textbf{Why homogeneous coordinates?}
\begin{itemize}
  \item Combine rotation and translation into a single matrix multiplication.
  \item Represent points at infinity with finite coordinates.
  \item Simplify composition of multiple transformations.
\end{itemize}

\separator

\subsubsection{Homogeneous Transformation Matrices (2D)}

\begin{defbox}[2D Homogeneous Transformation]
A rigid body transformation (rotation $\theta$ followed by translation $t = (x_t, y_t)^T$):
\[
T = \begin{pmatrix} R(\theta) & t \\ 0 & 1 \end{pmatrix} = \begin{pmatrix} \cos\theta & -\sin\theta & x_t \\ \sin\theta & \cos\theta & y_t \\ 0 & 0 & 1 \end{pmatrix}
\]
Applied to point $v$:
\[
\begin{pmatrix} v' \\ 1 \end{pmatrix} = T \begin{pmatrix} v \\ 1 \end{pmatrix} = \begin{pmatrix} R(\theta) v + t \\ 1 \end{pmatrix}
\]
\end{defbox}

\begin{thmbox}[Structure of Homogeneous Transform]
\[
T = \begin{pmatrix} \text{Rotation} & \text{Translation} \\ \text{Perspective} & \text{Scale} \end{pmatrix} = \begin{pmatrix} R & t \\ 0 & 1 \end{pmatrix}
\]
For rigid body transforms: perspective $= 0$, scale $= 1$.
\end{thmbox}

\begin{exbox}[Successive Transformations]
\textbf{Sequence:} Rotate by $\theta_1 \to$ Translate by $t_1 \to$ Rotate by $\theta_2 \to$ Translate by $t_2$

\textbf{Combined matrix:}
\[
T_{1,2} = \begin{pmatrix} R(\theta_2) & t_2 \\ 0 & 1 \end{pmatrix} \begin{pmatrix} R(\theta_1) & t_1 \\ 0 & 1 \end{pmatrix} = \begin{pmatrix} R(\theta_2)R(\theta_1) & R(\theta_2)t_1 + t_2 \\ 0 & 1 \end{pmatrix}
\]

\textbf{Key insight:} Transformations are applied right-to-left (matrix on the right acts first).
\end{exbox}

\separator

\subsubsection{Rotation About an Arbitrary Point (2D)}

\begin{thmbox}[Rotation About Point $p$]
To rotate about point $p$ (not the origin):
\begin{enumerate}
  \item Translate by $-p$ (move $p$ to origin)
  \item Rotate by $\theta$
  \item Translate by $+p$ (move back)
\end{enumerate}
\[
T = \begin{pmatrix} I & p \\ 0 & 1 \end{pmatrix} \begin{pmatrix} R(\theta) & 0 \\ 0 & 1 \end{pmatrix} \begin{pmatrix} I & -p \\ 0 & 1 \end{pmatrix}
\]
\end{thmbox}

\bigseparator

% =============================================================================
%                         3D TRANSFORMATIONS
% =============================================================================
\subsection{3D Transformations}

A rigid body in 3D requires:
\begin{itemize}
  \item \textbf{Position:} $(x, y, z) \in \mathbb{R}^3$ (3 parameters)
  \item \textbf{Orientation:} 3 parameters (various representations)
\end{itemize}

\subsubsection{3D Rotation Matrices}

\begin{defbox}[Elementary Rotation Matrices]
Rotations about the coordinate axes (right-hand rule, counter-clockwise):

\textbf{Roll} (about $x$-axis):
\[
R_x(\gamma) = \begin{pmatrix} 1 & 0 & 0 \\ 0 & \cos\gamma & -\sin\gamma \\ 0 & \sin\gamma & \cos\gamma \end{pmatrix}
\]

\textbf{Pitch} (about $y$-axis):
\[
R_y(\beta) = \begin{pmatrix} \cos\beta & 0 & \sin\beta \\ 0 & 1 & 0 \\ -\sin\beta & 0 & \cos\beta \end{pmatrix}
\]

\textbf{Yaw} (about $z$-axis):
\[
R_z(\alpha) = \begin{pmatrix} \cos\alpha & -\sin\alpha & 0 \\ \sin\alpha & \cos\alpha & 0 \\ 0 & 0 & 1 \end{pmatrix}
\]
\end{defbox}

\begin{tipbox}[Exam Tip: Memory Aid]
Each rotation matrix has 1 on the diagonal for the axis of rotation (that coordinate is unchanged). The $2\times 2$ block in the other positions is the standard 2D rotation matrix.
\end{tipbox}

\separator

\subsubsection{Euler Angles}

\begin{defbox}[Euler Angles]
Any 3D orientation can be achieved by \textbf{three successive rotations} about coordinate axes, where no two successive rotations are about the same axis.

\textbf{Common convention (ZYX / yaw-pitch-roll):}
\[
R(\alpha, \beta, \gamma) = R_z(\alpha) R_y(\beta) R_x(\gamma)
\]
Read right-to-left: roll $\gamma$ about $x$, then pitch $\beta$ about $y$, then yaw $\alpha$ about $z$.
\end{defbox}

\textbf{Valid Euler angle sequences:}
\begin{itemize}
  \item Proper Euler: $XYX, XZX, YXY, YZY, ZXZ, ZYZ$
  \item Tait-Bryan: $XYZ, XZY, YXZ, YZX, ZXY, ZYX$ (includes yaw-pitch-roll)
\end{itemize}

\begin{tipbox}[3D Rotations Do NOT Commute]
Unlike 2D, 3D rotations do not commute:
\[
R_z\left(\frac{\pi}{2}\right) R_y\left(\frac{\pi}{2}\right) \neq R_y\left(\frac{\pi}{2}\right) R_z\left(\frac{\pi}{2}\right)
\]
Order matters!
\end{tipbox}

\separator

\subsubsection{Properties of 3D Rotation Matrices}

\begin{thmbox}[Special Orthogonal Group $SO(3)$]
3D rotation matrices belong to $SO(3)$:
\[
SO(3) = \{R \in \mathbb{R}^{3\times 3} : R R^T = I \text{ and } \det(R) = 1\}
\]
\begin{itemize}
  \item Columns are mutually orthonormal: $r_i^T r_j = \delta_{ij}$
  \item Inverse: $R^{-1} = R^T$
  \item $\det(R) = +1$ (not $-1$, which would be a reflection)
\end{itemize}
\end{thmbox}

\separator

\subsubsection{Axis-Angle Representation}

\begin{defbox}[Rotation About Arbitrary Axis]
Any 3D rotation can be described by:
\begin{itemize}
  \item Unit vector $\mathbf{k} = (k_x, k_y, k_z)$ along the rotation axis
  \item Angle $\theta$ of rotation about that axis
\end{itemize}
The rotation matrix is computed by transforming to align with a coordinate axis:
\[
R_{\mathbf{k}}(\theta) = R_z(\alpha) R_y(\beta) R_z(\theta) R_y(-\beta) R_z(-\alpha)
\]
where $\alpha, \beta$ align $\mathbf{k}$ with the $z$-axis.
\end{defbox}

\separator

\subsubsection{Gimbal Lock}

\begin{defbox}[Gimbal Lock]
\term{Gimbal lock} occurs when two rotation axes become aligned, causing \textbf{loss of one degree of freedom}.

\textbf{Example:} When pitch $\beta = \pm 90^\circ$, yaw and roll rotations produce the same effect.
\end{defbox}

\begin{tipbox}[Gimbal Lock Problem]
\begin{itemize}
  \item Occurs with \emph{all} Euler angle representations (just at different configurations).
  \item Historical example: Apollo missions had to avoid certain orientations.
  \item Solution: Use quaternions!
\end{itemize}
\end{tipbox}

\separator

\subsubsection{Quaternions}

\begin{defbox}[Quaternion]
A \term{quaternion} is a 4-dimensional extension of complex numbers:
\[
q = a + b\mathbf{i} + c\mathbf{j} + d\mathbf{k} = (a, \mathbf{v}) \quad \text{where } \mathbf{v} = (b, c, d)
\]
With multiplication rules:
\[
\mathbf{i}^2 = \mathbf{j}^2 = \mathbf{k}^2 = -1, \quad \mathbf{ij} = \mathbf{k}, \quad \mathbf{jk} = \mathbf{i}, \quad \mathbf{ki} = \mathbf{j}
\]
\end{defbox}

\begin{defbox}[Unit Quaternion for Rotation]
A \term{unit quaternion} ($\|q\| = 1$) represents rotation by angle $\theta$ about unit axis $\mathbf{n}$:
\[
q = \cos\frac{\theta}{2} + \sin\frac{\theta}{2}(n_x\mathbf{i} + n_y\mathbf{j} + n_z\mathbf{k}) = e^{\frac{\theta}{2}\mathbf{n}}
\]
\end{defbox}

\begin{thmbox}[Rotating a Point with Quaternions]
To rotate point $(x, y, z)$ by quaternion $q$:
\begin{enumerate}
  \item Represent point as pure quaternion: $v = x\mathbf{i} + y\mathbf{j} + z\mathbf{k}$
  \item Compute: $v' = q \cdot v \cdot q^*$ \quad (quaternion ``sandwich'')
  \item Extract rotated point from $v' = x'\mathbf{i} + y'\mathbf{j} + z'\mathbf{k}$
\end{enumerate}
where $q^* = a - b\mathbf{i} - c\mathbf{j} - d\mathbf{k}$ is the conjugate.
\end{thmbox}

\begin{keybox}[Advantages of Quaternions]
\begin{center}
\renewcommand{\arraystretch}{1.3}
\begin{tabular}{@{} l p{9cm} @{}}
\toprule
\textbf{Property} & \textbf{Benefit} \\
\midrule
Compact & 4 parameters vs.\ 9 for rotation matrix \\
No gimbal lock & Avoids singularities entirely \\
Easy renormalization & Just normalize to unit length \\
Stable interpolation & SLERP gives smooth rotation interpolation \\
Numerically stable & Avoids accumulating errors from matrix multiplication \\
\bottomrule
\end{tabular}
\end{center}
\textbf{Disadvantage:} Not intuitive; must convert to matrix to visualize.
\end{keybox}

\begin{defbox}[SLERP: Spherical Linear Interpolation]
Interpolate between quaternions $q_1$ and $q_2$:
\[
\text{SLERP}(q_1, q_2, t) = \frac{\sin((1-t)\phi)}{\sin\phi} q_1 + \frac{\sin(t\phi)}{\sin\phi} q_2
\]
where $\phi$ is the angle between $q_1$ and $q_2$. Always produces a unit quaternion.
\end{defbox}

\separator

\subsubsection{3D Homogeneous Transformations}

\begin{defbox}[3D Homogeneous Transformation Matrix]
Rotation $R$ followed by translation $t = (x_t, y_t, z_t)^T$:
\[
T = \begin{pmatrix} R & t \\ 0 & 1 \end{pmatrix} \in \mathbb{R}^{4\times 4}
\]
Applied to point $v$:
\[
\begin{pmatrix} v' \\ 1 \end{pmatrix} = T \begin{pmatrix} v \\ 1 \end{pmatrix}
\]
\end{defbox}

\bigseparator

% =============================================================================
%                              KINEMATICS
% =============================================================================
\subsection{Kinematics}

\begin{defbox}[Kinematics]
\term{Kinematics} is the study of motion and configurations of a system---the ``geometry of the system''---without considering forces.

Key concepts:
\begin{itemize}
  \item \term{Link}: A rigid body in a chain of rigid bodies.
  \item \term{Joint}: Connects two links and constrains their relative motion.
  \item \term{End-effector}: The ``hand'' or tool at the end of the kinematic chain.
\end{itemize}
\end{defbox}

\separator

\subsubsection{Forward Kinematics}

\begin{defbox}[Forward Kinematics (FK)]
Given joint angles/positions $\mathbf{q} = (q_1, \ldots, q_n)$, compute the position and orientation of the end-effector in the world frame:
\[
\mathbf{x} = FK(\mathbf{q})
\]
\end{defbox}

\begin{thmbox}[FK via Transformation Matrices]
For a chain of $n$ links, each with transformation $T_i$ from link $i$ to link $i-1$:
\[
T_{\text{total}} = T_1 T_2 \cdots T_n T_{\text{ee}}
\]
where $T_{\text{ee}}$ is the end-effector offset from the last joint.

The end-effector position in the world frame:
\[
\begin{pmatrix} x \\ y \\ 1 \end{pmatrix} = T_{\text{total}} \begin{pmatrix} 0 \\ 0 \\ 1 \end{pmatrix}
\]
\end{thmbox}

\begin{exbox}[2D 3-Link Planar Arm]
For a 3-link arm with link lengths $a_1, a_2, a_3$ and joint angles $\theta_1, \theta_2, \theta_3$:

\textbf{Transformation matrices:}
\[
T_i = \begin{pmatrix} \cos\theta_i & -\sin\theta_i & a_{i-1} \\ \sin\theta_i & \cos\theta_i & 0 \\ 0 & 0 & 1 \end{pmatrix}
\]

\textbf{End-effector position:}
\begin{align*}
x &= a_1\cos\theta_1 + a_2\cos(\theta_1+\theta_2) + a_3\cos(\theta_1+\theta_2+\theta_3) \\
y &= a_1\sin\theta_1 + a_2\sin(\theta_1+\theta_2) + a_3\sin(\theta_1+\theta_2+\theta_3)
\end{align*}
\end{exbox}

\separator

\subsubsection{Inverse Kinematics}

\begin{defbox}[Inverse Kinematics (IK)]
Given a desired end-effector position/orientation $\mathbf{x}_{\text{target}}$, find joint angles $\mathbf{q}$ such that:
\[
FK(\mathbf{q}) = \mathbf{x}_{\text{target}}
\]
\end{defbox}

\begin{tipbox}[IK Solution Cases]
Let $n$ = number of joint DOFs, $m$ = number of task DOFs.
\begin{center}
\renewcommand{\arraystretch}{1.3}
\begin{tabular}{@{} c l @{}}
\toprule
\textbf{Condition} & \textbf{Solutions} \\
\midrule
$n < m$ & No solution (overconstrained) \\
$n = m$ & Finite number of solutions (if reachable) \\
$n > m$ & Infinite solutions (robot is \term{redundant}) \\
\bottomrule
\end{tabular}
\end{center}
\end{tipbox}

\begin{exbox}[Analytic IK for 2-Link Planar Arm]
\textbf{Given:} Target $(x, y)$, link lengths $a_1, a_2$.

\textbf{Forward kinematics:}
\begin{align*}
x &= a_1\cos\theta_1 + a_2\cos(\theta_1 + \theta_2) \\
y &= a_1\sin\theta_1 + a_2\sin(\theta_1 + \theta_2)
\end{align*}

\textbf{Solution for $\theta_2$:} Square and add the FK equations:
\[
x^2 + y^2 = a_1^2 + a_2^2 + 2a_1 a_2 \cos\theta_2
\]
\[
\cos\theta_2 = \frac{x^2 + y^2 - a_1^2 - a_2^2}{2a_1 a_2}
\]
\[
\sin\theta_2 = \pm\sqrt{1 - \cos^2\theta_2}
\]

\textbf{Solution for $\theta_1$:} Substitute back:
\[
\cos\theta_1 = \frac{1}{x^2+y^2}\left[x(a_1 + a_2\cos\theta_2) \pm y \cdot a_2\sin\theta_2\right]
\]
\[
\sin\theta_1 = \frac{1}{x^2+y^2}\left[y(a_1 + a_2\cos\theta_2) \mp x \cdot a_2\sin\theta_2\right]
\]

\textbf{Two solutions:} ``Elbow up'' and ``elbow down'' configurations.
\end{exbox}

\begin{defbox}[Singularity]
A \term{singularity} occurs when the robot loses one or more degrees of freedom---only one IK solution exists instead of two (or the Jacobian loses rank).

\textbf{Examples:}
\begin{itemize}
  \item Arm fully extended: $\theta_2 = 0$
  \item Arm fully folded: $\theta_2 = \pi$
\end{itemize}
\end{defbox}

\separator

\subsubsection{The Manipulator Jacobian}

\begin{defbox}[Jacobian Matrix]
The \term{Jacobian} $J$ relates joint velocities to end-effector velocities:
\[
\dot{\mathbf{x}} = J(\mathbf{q}) \cdot \dot{\mathbf{q}}
\]
where:
\[
J = \frac{\partial FK}{\partial \mathbf{q}} = \begin{pmatrix} \frac{\partial x}{\partial q_1} & \cdots & \frac{\partial x}{\partial q_n} \\ \frac{\partial y}{\partial q_1} & \cdots & \frac{\partial y}{\partial q_n} \\ \vdots & & \vdots \end{pmatrix}
\]
\end{defbox}

\begin{thmbox}[Numerical IK via Jacobian]
To solve IK iteratively:
\[
\Delta \mathbf{q} = J^{-1} \Delta \mathbf{x}
\]
where $\Delta \mathbf{x} = \mathbf{x}_{\text{target}} - FK(\mathbf{q})$.

\textbf{Gradient descent approach:}
\begin{enumerate}
  \item Compute error: $\Delta \mathbf{x} = \mathbf{x}_{\text{target}} - FK(\mathbf{q})$
  \item Compute Jacobian $J(\mathbf{q})$
  \item Update: $\mathbf{q} \leftarrow \mathbf{q} + \alpha J^{-1} \Delta \mathbf{x}$
  \item Repeat until $\|\Delta \mathbf{x}\| < \epsilon$
\end{enumerate}

For non-square $J$, use pseudoinverse $J^+$ or transpose $J^T$.
\end{thmbox}

\begin{tipbox}[Jacobian and Singularities]
At a \textbf{singularity}, the Jacobian loses rank (rows become linearly dependent). This means:
\begin{itemize}
  \item Some end-effector velocities are impossible.
  \item $J^{-1}$ is undefined or numerically unstable.
  \item Use \textbf{damped least squares} to handle near-singularities.
\end{itemize}
\end{tipbox}

\bigseparator

\subsection{Summary Tables}

\begin{center}
\renewcommand{\arraystretch}{1.3}
\begin{tabular}{@{} l c c c @{}}
\toprule
\textbf{Rotation Representation} & \textbf{Parameters} & \textbf{Gimbal Lock?} & \textbf{Interpolation} \\
\midrule
Rotation matrix & 9 (with 6 constraints) & No & Difficult \\
Euler angles & 3 & Yes & Linear (discontinuous) \\
Axis-angle & 4 (with 1 constraint) & At $\theta=0$ & Difficult \\
Quaternion & 4 (with 1 constraint) & No & SLERP (smooth) \\
\bottomrule
\end{tabular}
\end{center}

\bigseparator

\subsection{Worked Examples}

\begin{exbox}[Example 1: 2D Transformation Sequence]
\textbf{Problem:} A robot at origin is rotated by $90^\circ$ then translated by $(3, 2)$. Find the transformation matrix and the final position of point $(1, 0)$.

\textbf{Solution:}
\[
T = \begin{pmatrix} \cos 90^\circ & -\sin 90^\circ & 3 \\ \sin 90^\circ & \cos 90^\circ & 2 \\ 0 & 0 & 1 \end{pmatrix} = \begin{pmatrix} 0 & -1 & 3 \\ 1 & 0 & 2 \\ 0 & 0 & 1 \end{pmatrix}
\]
\[
\begin{pmatrix} x' \\ y' \\ 1 \end{pmatrix} = \begin{pmatrix} 0 & -1 & 3 \\ 1 & 0 & 2 \\ 0 & 0 & 1 \end{pmatrix} \begin{pmatrix} 1 \\ 0 \\ 1 \end{pmatrix} = \begin{pmatrix} 3 \\ 3 \\ 1 \end{pmatrix}
\]
Final position: $(3, 3)$.
\end{exbox}

\begin{exbox}[Example 2: 2-Link IK]
\textbf{Problem:} A 2-link arm has $a_1 = a_2 = 1$. Find $\theta_1, \theta_2$ for end-effector at $(1, 1)$.

\textbf{Solution:}
\[
\cos\theta_2 = \frac{1^2 + 1^2 - 1 - 1}{2(1)(1)} = \frac{0}{2} = 0 \implies \theta_2 = \pm 90^\circ
\]

For $\theta_2 = 90^\circ$: $\sin\theta_2 = 1$
\[
\cos\theta_1 = \frac{1(1+0) + 1(1)}{2} = 1, \quad \sin\theta_1 = \frac{1(1+0) - 1(1)}{2} = 0
\]
$\theta_1 = 0^\circ$

Check: $x = \cos 0^\circ + \cos 90^\circ = 1$, $y = \sin 0^\circ + \sin 90^\circ = 1$. \checkmark
\end{exbox}

\begin{exbox}[Example 3: Quaternion Rotation]
\textbf{Problem:} Express a $180^\circ$ rotation about the $z$-axis as a quaternion.

\textbf{Solution:}
\[
q = \cos\frac{180^\circ}{2} + \sin\frac{180^\circ}{2}(0\mathbf{i} + 0\mathbf{j} + 1\mathbf{k}) = \cos 90^\circ + \sin 90^\circ \cdot \mathbf{k} = 0 + 1\mathbf{k} = \mathbf{k}
\]
Or in $(a, b, c, d)$ form: $q = (0, 0, 0, 1)$.
\end{exbox}
